\chapter{Introdução}
\label{cap:01}

Os jogos eletrônicos é um dos setores mais relevantes da indústria do entretenimento,
movimentando bilhões de dólares e engajando milhões de jogadores ao redor do mundo, 
a indústria de jogos eletrônicos gera mais que a indústrial músical e do cinema somadas,
segundo o newzoo cerca de 183 bilhões foram movimentados no mercado de jogos eletrônicos
para consoles, computadores e celulares \cite{Newzoo, Statista2025GamesMarket}, além de sua importância cultural
e econômica, os jogos também são um ótimo fator para estudos acadêmicos, inclusive no
estudo de desenvolvimento de software, aprender lógica de programação, uso de diferentes tecnologias e à 
aplicação de distintos paradigmas de programação.
Fazer essa comparação das escolhas torna-se importante não apenas para a prática
acadêmica, mas também para orientar iniciantes e profissionais na área de desenvolvimento de jogos e 
descobrir qual o impacto da escolha da arquitetura e do paradigma na produtividade do desenvolvedor e na 
performance final de um jogo.

Com base nessas informações, será feito um teste de hipóteses analisando e comparando diferentes tecnologias de desenvolvimento de software,
utilizando o jogo \textit{Tetris} para ser feita essa análise, foi escolhido \textit{Tetris} devido a sua estrutura ser simples e trivial.


\section{Objetivos}

\subsection{Objetivo Geral}

O objetivo geral do trabalho é comparar quatro formas diferentes de implementação do jogo \textit{Tetris}, por meio
das tecnologias \textit{Godot}, Java , C++ e Python.

\subsection{Objetivos Específicos}
\begin{itemize}
	\item Desenvolver Tetris no motor Gráfico \textit{Godot};
	\item Desenvolver Tetris com a linguagem Java utilizando a biblioteca LibGDX;
	\item Desenvolver Tetris com a linguagem c++ utilizando a biblioteca Raylib;
	\item Desenvolver Tetris com a linguagem Python utilizando a biblioteca PyGame;
	\item Coletar dados de tempo de desenvolvimento e perfomance dos jogos de Tetris;
	\item Comparar dados de tempo de desenvolvimento e perfomance dos jogos de Tetris;
\end{itemize}